% Ad Familiares 6.8 --- headnote
% ad-familiares-06-08, early December 46 BC, Rome

\textit{Cicero to Aulus Caecina, written at Rome at the
beginning of December 46 BC (Perseus: \textit{Romae in.~m.~Dec.
a.~708 (46)}; works.yaml carries mid-December at month
precision). The letter closes the Caecina cluster
(\textit{Fam.}~6.5--8); for Caecina see the headnote to
\textit{Fam.}~6.6. The piece is the practical companion to the
substantive \textit{Fam.}~6.6: where that letter argued the
political reading, this one reports a concrete success in
negotiation and disposes of Caecina's logistical question.}

\textit{Three short paragraphs. In \S~1 Cicero has, through
Balbus and Oppius --- Caesar's regular men of business in his
absence --- got Caecina's deadline of 1~January waived, with a
guarantee that he may remain in Sicily as long as he wishes
without incurring Caesar's displeasure. (Cicero notes that what
Balbus and Oppius transact in Caesar's absence is taken by
Caesar as ratified; this is the practical point of the cluster.)
In \S~2 he answers a letter that had crossed his own --- Caecina
had asked whether he should stay in Sicily or set out to settle
his remaining affairs in Asia. Cicero's reading is unambiguous:
stay; the nearness of Sicily helps both the petition and the
speed of the eventual return. In \S~3 he commends Caecina to
Titus Furfanius Postumus, the new governor of Sicily, and his
legates --- all friends of Cicero's, all just arrived from
Mutina --- and encloses a letter of recommendation to Furfanius,
appended below in the manuscript tradition (it is not preserved
in our text). This is the form of mediation Caecina had asked
for in \textit{Fam.}~6.7: the load on Cicero, the moves coming
from Cicero, the issue brought through by Cicero. The reply
delivers.}
